% =============================================================================
% CoT-Faith figures — LaTeX include snippets (paste into cot_faith.tex)
% =============================================================================
%
% Required packages in preamble:
%   \usepackage{graphicx}
%   \usepackage{tikz}
%   \usetikzlibrary{arrows.meta, positioning}
%
% All PDFs are vector (fonts embedded, Type 42), rendered at 300 DPI.

% =============================================================================
% Fig 1 — CoT-Faith workflow overview (TikZ, MANUAL polish)
% =============================================================================
% See figures/fig1_workflow.tex for the raw TikZ. Recommended: paste inline
% at the start of Section 1 (Introduction) as the hero/teaser figure.
%
%   % =============================================================
% Fig 1: CoT-Faith workflow overview (TikZ skeleton — POLISH BY HAND)
% =============================================================
% Suggested placement: after \begin{document}, floated at [t] on p.1.
% Compile with pdflatex (loads tikz). Adjust node positions/colours
% to match the paper's colour scheme.
%
\begin{figure}[t]
\centering
\begin{tikzpicture}[
    node distance=0.7cm,
    box/.style={rectangle, rounded corners=2pt, draw=black, thick,
                minimum width=2.4cm, minimum height=0.9cm,
                align=center, font=\small},
    inbox/.style={box, fill=blue!8},
    modbox/.style={box, fill=green!12},
    probe/.style={box, fill=orange!12, minimum width=3.0cm},
    findbox/.style={box, fill=red!7, minimum width=3.4cm, minimum height=0.8cm},
    arr/.style={-{Latex[length=2mm]}, thick, shorten <=1pt, shorten >=1pt},
]

% Row 1: inputs
\node[inbox]  (img) at (0, 0)  {LIBERO image\\ ({\scriptsize$H\times W\times3$})};
\node[inbox, right=0.5cm of img]  (instr) {instruction\\ ({\scriptsize"pick up X"})};
\node[inbox, right=0.5cm of instr] (cot)   {\emph{GT} CoT\\ ({\scriptsize9-tag trace})};

% Row 2: model
\node[modbox, below=0.9cm of instr, minimum width=6.4cm, minimum height=1.0cm]
      (model) {CoT-VLA (ECoT-family or ablation variant)};

% Row 3: three probes
\node[probe, below=1.1cm of model.south west, xshift=0.6cm] (p1)
      {\textbf{Probe 1}: segment\\ attention (4 buckets)};
\node[probe, below=1.1cm of model] (p2)
      {\textbf{Probe 2}: causal\\ CoT edit ($\Delta_\infty$)};
\node[probe, below=1.1cm of model.south east, xshift=-0.6cm] (p3)
      {\textbf{Probe 3}: attention\\ $\to$ error AUROC};

% Row 4: findings
\node[findbox, below=0.9cm of p1] (f1)
      {attn shape = arch-\\determined (Fig 2)};
\node[findbox, below=0.9cm of p2] (f2)
      {reasoning target\\ = causal (Fig 3-4)};
\node[findbox, below=0.9cm of p3] (f3)
      {attn $\ne$ error\\ AUROC $\approx 0.5$};

% Arrows: inputs -> model
\draw[arr] (img)   -- ($(model.north)+(-1.8,0)$);
\draw[arr] (instr) -- (model.north);
\draw[arr] (cot)   -- ($(model.north)+(1.8,0)$);

% Model -> probes
\draw[arr] (model.south) -- (p1.north);
\draw[arr] (model.south) -- (p2.north);
\draw[arr] (model.south) -- (p3.north);

% Probes -> findings
\draw[arr] (p1) -- (f1);
\draw[arr] (p2) -- (f2);
\draw[arr] (p3) -- (f3);

\end{tikzpicture}
\caption{\textbf{CoT-Faith methodology}. Each LIBERO first-step
sample $(\mathrm{img}, \mathrm{instr}, \mathrm{cot}_\text{GT})$ is fed
to a CoT-VLA with the standard 9-tag ECoT prompt. Three probes measure
different aspects of chain-of-thought faithfulness: segment-decomposed
action-token attention (Probe 1); Δaction under causal edits along ten
semantic axes with a null selfsplice control (Probe 2); AUROC of
attention as a per-sample action-error predictor (Probe 3). Across
11 models we find that attention on the CoT segment is
architecture-determined, while causal effect requires reasoning-target
training — attention does not proxy causation.}
\label{fig:workflow}
\end{figure}


% =============================================================================
% Fig 2 — Action-token attention across 11 models
% =============================================================================
\begin{figure}[t]
    \centering
    \includegraphics[width=0.95\textwidth]{figures/fig2_attention_distribution.pdf}
    \caption{\textbf{Action-token attention distribution across 11 models}.
    Bars show mean attention mass from action-token positions to each
    of four sequence segments (visual, instruction, CoT, action-prev),
    averaged over LLM layers 0-3, $N{=}100$ LIBERO first-step samples, and
    up to 3 seeds. All CoT-VLAs (left of dashed line) cluster at
    $(\text{visual}, \text{instr}, \text{CoT}, \text{prev}) \approx
    (0.29, 0.30, 0.34, 0.07)$; non-CoT baselines (right) cluster at
    $(0.52, 0.42, 0.00, 0.05)$. Critically, our no-CoT model —
    fine-tuned on the same LIBERO-90 data as our r=32 model but with
    all reasoning tags stripped from the training loss — still shows
    the CoT-VLA attention shape when fed a CoT-formatted prompt.}
    \label{fig:attention}
\end{figure}

% =============================================================================
% Fig 3 — Causal-edit faithful-rate heatmap
% =============================================================================
\begin{figure}[t]
    \centering
    \includegraphics[width=0.95\textwidth]{figures/fig3_edit_heatmap.pdf}
    \caption{\textbf{Causal-edit faithful rate across 8 CoT-VLAs $\times$ 10
    edit families}. Each cell is the fraction of $N{=}100$ LIBERO
    samples where the CoT edit produced $\Delta_\infty > 0.05$ in the
    greedy-decoded action. Two anchor columns: \texttt{selfsplice} (identity
    replacement — validates the metric, must be $\approx 0$) and
    \texttt{cross\_task\_swap} (replace whole CoT with unrelated
    sample's — upper bound). \textbf{Row highlights:}
    \textcolor[HTML]{EE6677}{\textbf{no-CoT}} collapses to
    0.13-0.24 on every semantic edit despite identical base and data,
    while \textcolor[HTML]{228833}{\textbf{ECoT-bridge}} sits at
    0.67-0.97 on the same edits. Missing cells (\texttt{---}) are edits
    not yet re-run for our r=32 model.}
    \label{fig:edit_heatmap}
\end{figure}

% =============================================================================
% Fig 4 — Attention-vs-causation dissociation
% =============================================================================
\begin{figure}[t]
    \centering
    \includegraphics[width=0.95\textwidth]{figures/fig4_dissociation.pdf}
    \caption{\textbf{Attention-causation dissociation}.
    \textbf{(a)} Action-token attention on the CoT segment for 7
    trained CoT-VLA variants: all cluster at $\sim 0.34$, with the
    no-CoT ablation (red) at 0.35 — indistinguishable from
    CoT-trained models. \textbf{(b)} Same models, average
    faithful rate across seven shared edit families (subject swap,
    direction flip, gripper flip, verb swap, negation,
    adversarial-plausible, cross-task swap): no-CoT collapses to
    0.15, while ECoT-bridge (green) hits 0.87. \textbf{Attention
    allocation does not predict causal effect}: the same 34\%
    attention on CoT can drive the action decode (bridge) or be
    entirely decorative (no-CoT).}
    \label{fig:dissociation}
\end{figure}

% =============================================================================
% Fig 5 — Prompt-format ablation
% =============================================================================
\begin{figure}[t]
    \centering
    \includegraphics[width=0.6\textwidth]{figures/fig5_prompt_ablation.pdf}
    \caption{\textbf{Prompt-format ablation on ECoT-bridge}.
    Fraction of $N{=}100$ LIBERO samples where truncating the CoT
    text causes $\Delta_\infty > 0.05$ in the action. Removing
    \emph{any} portion of the 9-tag CoT — even leaving just the
    single-word \texttt{TASK:} tag — causes $\geq 96\%$ of samples
    to change action. Shuffling tag order alone triggers 95\%.
    The model conditions strongly on the full structured target;
    but this is a property of the ECoT-bridge \emph{training regime}
    (Fig~\ref{fig:edit_heatmap} shows models without CoT training
    do not exhibit this sensitivity).}
    \label{fig:prompt}
\end{figure}

% =============================================================================
% Fig 6 — Bridge vs LIBERO CoT-training comparison
% =============================================================================
\begin{figure}[t]
    \centering
    \includegraphics[width=0.55\textwidth]{figures/fig6_bridge_vs_libero.pdf}
    \caption{\textbf{Bridge-trained vs LIBERO-fine-tuned CoT-VLA
    faithful rate}. On three shared edit families, ECoT-bridge
    (green) achieves 0.67-0.97 faithful rate — approximately
    $2\times$ our r=32 LIBERO fine-tune (blue, 0.24-0.77) despite
    identical architecture, prompt format, and near-identical
    action-token attention distributions (Fig~\ref{fig:attention}).
    CoT-Faith is thus governed by training-data domain, not
    architecture.}
    \label{fig:bridge_vs_libero}
\end{figure}
